\section{Conclusion}
\label{sec:conclusion}

% ⚠️ 本节的【一句话核心结论】必须与以下三处严格一致：
%    - report/sections/00-abstract.tex
%    - ../README.md 第 2 节
%    - ../slides/outline.md 的 P1 与 P13
%
%    这是课程硬指标之一（1 句核心结论）。请在提交前逐一核对。

\subsection*{One-sentence claim}

% 填写说明（写完请删除本段注释）：
%   要求是一【句】结论，不是一段。
%   必须是我们从文献中【推导出的判断】，而非【本文综述了什么】。
%
%   反例（不合格）:
%     "This review surveys recent methods of harness self-improvement."
%   正例（合格）:
%     "Under condition X, the gains reported in harness self-improvement are
%      attributable to cause Y, but that explanation fails when condition Z."
%
% ⚠️ colorbox 内部不能有空行，否则 \par 会打断分组。

\begin{quote}
\large
\colorbox{yellow!15}{%
  \parbox{\dimexpr\linewidth-2\fboxsep\relax}{%
    \todo{Write the one-sentence core claim here.}%
  }%
}
\end{quote}

\subsection{What we established}

\todo{Three or four points, one line each, each pointing to the corresponding
section and evidence item.}

\subsection{What remains open}

\todo{Two or three points, strictly matching Section~\ref{sec:open}.}

\subsection{Limitations of this review}

\todo{Required. At minimum acknowledge: the corpus is small (n = \todo{N}),
so claims about distribution must be cautious; most of the corpus is
preprints whose venue status is not fully verified; we reproduced no
experiments, so all conclusions rest on reading and comparison of the
literature.}

\subsection{Outlook}

\todo{Two or three sentences on the most promising next step. The empty cells
recorded in taxonomy.md are a natural place to start: formal guarantees, and
empirical work on co-evolution.}

% ----------------------------------------------------------------------
% 相关中文文件（不参与编译，仅供对照）:
%   ../synthesis/evidence.md          本节内容来源
%   ../synthesis/open-questions.md    Section 5 内容来源
%   ../synthesis/taxonomy.md          Section 3 内容来源
% ----------------------------------------------------------------------
